% Prozentrechnung · Umwandeln · Prüfungs-Fokus MSA – bank/prozentrechnung/e1.jsonl (zusammenbau v0.6, Prüfungs-Fokus)
\einheitenkopf[e1]{Einheit 1 · Prozente als Anteile}
\zweigzeile{ab Kl. 6 (am Gymnasium ab Kl. 5) · P10 ×6}
\setcounter{aufgabe}{0}

% Anlauf (ohne Prüfkennung)
\begin{aufgabe}{Umwandeln – Anlauf}
\begin{teile}
\teil Wie viel Prozent sind grau?

\streifenfeld[2]{50}{$0\,\%$}{$100\,\%$}
\teil $\frac{37}{100}$ – wie viel Prozent? \leerfeld[\%]
\end{teile}
\end{aufgabe}

\begin{aufgabe}{Umwandeln – Prüfungsaufgaben}
\begin{teile}
\teil $5\,\%$ der Pakete kommen zu spät – welche Aussage passt? Kreuze an. \hfill (P20)\\ \kreuz{$5$ von $10$ Paketen kommen zu spät.}\\ \kreuz{Jedes 5. Paket kommt zu spät.}\\ \kreuz{$5$ von $100$ Paketen kommen zu spät.}
\teil $8\,\%$ der Gäste bestellen Tee – welche Aussage passt? Kreuze an. \hfill (P20)\\ \kreuz{Jeder 8. Gast bestellt Tee.}\\ \kreuz{$8$ von $100$ Gästen bestellen Tee.}\\ \kreuz{$8$ von $10$ Gästen bestellen Tee.}
\teil Jeder vierte Besucher kauft ein Programmheft – wie viel Prozent? Kreuze an. \hfill (P22)\\ \kreuz{$4\,\%$}\\ \kreuz{$25\,\%$}\\ \kreuz{$40\,\%$}\\ \kreuz{$75\,\%$}
\teil Jedes zwanzigste Ei einer Lieferung ist zerbrochen – wie viel Prozent? Kreuze an. \hfill (P22)\\ \kreuz{$2\,\%$}\\ \kreuz{$5\,\%$}\\ \kreuz{$20\,\%$}\\ \kreuz{$80\,\%$}
\end{teile}
\end{aufgabe}

\begin{aufgabe}{Umwandeln – Prüfungsaufgaben – weiter}
\begin{teile}
\teil Schulweg der Kinder einer Schule. Genau eine Aussage ist falsch. Kreuze sie an und schreibe sie richtig. \hfill (P23)\\ \kreuz{Aussage 1: Ein Viertel der Kinder fährt mit dem Bus.}\\ \kreuz{Aussage 2: Die Hälfte der Kinder fährt mit dem Rad.}

\sachtabelle{lr}{Schulweg & Kinder}{Bus & 150\\ Rad & 200\\ zu Fuß & 180\\ Auto & 70\\ gesamt & 600}
\teil Mitglieder eines Sportvereins. Genau eine Aussage ist falsch. Kreuze sie an und schreibe sie richtig. \hfill (P23)\\ \kreuz{Aussage 1: Jedes fünfte Mitglied ist ein Kind.}\\ \kreuz{Aussage 2: Die Hälfte der Mitglieder sind Erwachsene.}

\sachtabelle{lr}{Gruppe & Mitglieder}{Kinder & 600\\ Jugendliche & 480\\ Erwachsene & 1\,200\\ Senioren & 120\\ gesamt & 2\,400}
\teil Freizeit: Anteil der Befragten in der Stadt und auf dem Land. Genau eine Aussage ist falsch. Kreuze sie an und schreibe sie richtig. \hfill (P19)\\ \kreuz{Aussage 1: Ein Viertel der Befragten auf dem Land geht in die Bibliothek.}\\ \kreuz{Aussage 2: Mehr als die Hälfte der Befragten auf dem Land treibt Sport im Verein.}

\sachtabelle{lcc}{Freizeit & Stadt & Land}{Sportverein & 40\,\% & 55\,\%\\ Kino & 12\,\% & 5\,\%\\ Bibliothek & 25\,\% & 10\,\%\\ Schwimmbad & 60\,\% & 48\,\%}
\teil Frühstück zu Hause: Anteil der Kinder an der Grundschule und an der Oberschule. Genau eine Aussage ist falsch. Kreuze sie an und schreibe sie richtig. \hfill (P19)\\ \kreuz{Aussage 1: Drei von fünf Oberschülern frühstücken immer zu Hause.}\\ \kreuz{Aussage 2: Jeder zehnte Grundschüler frühstückt nie zu Hause.}

\sachtabelle{lcc}{Frühstück & Grundschule & Oberschule}{immer & 80\,\% & 60\,\%\\ manchmal & 15\,\% & 30\,\%\\ nie & 5\,\% & 10\,\%}
\steil Wie oft Jugendliche Sport treiben (Anteile). Jemand behauptet: „Mehr als die Hälfte der Jugendlichen treibt mindestens zweimal pro Woche Sport.“ Stimmt das? Begründe mit den Werten der Tabelle. \hfill (P17*)

\sachtabelle{lr}{Sport & Anteil}{nie & 18{,}5\,\%\\ einmal pro Woche & 32{,}0\,\%\\ zweimal pro Woche & 27{,}5\,\%\\ dreimal oder öfter & 22{,}0\,\%}
\steil Haushalte nach Zahl der Autos (Anteile). Jemand behauptet: „Mehr als $70\,\%$ der Haushalte haben mindestens ein Auto.“ Stimmt das? Begründe mit den Werten der Tabelle. \hfill (P17*)

\sachtabelle{lr}{Autos & Anteil}{kein Auto & 28{,}6\,\%\\ ein Auto & 45{,}3\,\%\\ zwei Autos & 21{,}8\,\%\\ drei oder mehr & 4{,}3\,\%}
\end{teile}
\end{aufgabe}

\begin{aufgabe}{Umwandeln – Prüfungsaufgaben – weiter}
\begin{teile}
\teil Strom einer Gemeinde nach Energiequelle – trage den fehlenden Anteil ein. \hfill (P21)

\sachtabelle{lr}{Energiequelle & Anteil}{Wind & 38\,\%\\ Sonne & 24\,\%\\ Biogas & \leerzelle\\ Wasser & 9\,\%\\ Sonstiges & 3\,\%\\ gesamt & 100\,\%}
\teil Fläche eines Landkreises nach Nutzung – trage den fehlenden Anteil ein. \hfill (P21)

\sachtabelle{lr}{Nutzung & Anteil}{Wald & 34\,\%\\ Acker & 41\,\%\\ Siedlung & \leerzelle\\ Wasser & 4\,\%\\ Sonstiges & 2\,\%\\ gesamt & 100\,\%}
\end{teile}
\end{aufgabe}

\medskip
\uebersichtskasten{Prozent heißt Hundertstel: 1 \% = 1/100 = 0,01. \\ Bruch → Prozent: auf Nenner 100 erweitern oder Dezimalzahl mit Komma zwei Stellen nach rechts. \\ 3/4 = 75/100 = 75 \% \quad 0,08 = 8 \% \quad „jeder Fünfte“ = 1/5 = 20 \% \\ Alle Anteile zusammen sind 100 \%.}
